% ============================================================
% §2.1  Novelty as agent-relative mismatch
% ============================================================

\section{Perspectives on Novelty}
\label{sec:novelty_as_mismatch}
\label{sec:novelty_perspectives}

The first step in organizing work on open-world adaptation is to clarify what is meant by novelty.
Across psychology, neuroscience, and artificial intelligence, novelty is best understood not as an intrinsic property of an object or event, but as a relation between the world and an agent's current memory, expectations, and models.
The same object, outcome, or environmental change may be routine for one agent and disruptive for another.
For this dissertation, this agent-relative view is central: novelty matters because it exposes a mismatch between what the world presents and what the agent can currently represent, predict, or do.

A useful starting point is the distinction between novelty and surprise.
\citet{barto2013novelty} argue that the two terms name different comparisons.
Novelty concerns whether something has been encountered or represented before; surprise concerns whether an expectation has been violated.
A familiar door that remains locked after an unlock action is surprising, because the outcome violates a prediction, but the door and lock are not novel categories to an agent that already represents them.
Conversely, an exploratory agent may encounter an object type it has never represented before without being strongly surprised if unfamiliar objects were expected during exploration.
The two signals often coincide, but they place different burdens on an agent.
Surprise calls for better prediction.
Novelty may require a new representation.
\label{sec:novelty_surprise}

Neuroscience gives this relational point a more mechanistic form.
Novel events engage brain systems involved in attention and memory, helping the agent orient toward unexpected information and encode it for later use~\citep{ranganath2003neural}.
Predictive-processing accounts make a related point: perception depends on predictions, and learning is driven in part by the mismatch between what is expected and what is observed~\citep{clark2013predictive,friston2010free}.
For this chapter, the important point is that mismatch is not all of one kind.
An agent may only need to refine a prediction within a representation it already has, or it may need to form a representation for something that was not previously expressible.
\label{sec:novelty_neuro}

Artificial intelligence reaches the same view from the problem of bounded agency.
The big world hypothesis argues that a realistic agent is far smaller than the world it inhabits and therefore cannot maintain a complete model of that world~\citep{javed2024bigworld}.
If the agent's model is necessarily partial, then encountering its limits is not an exceptional malfunction.
It is a normal condition of acting in a world whose structure exceeds the agent's representation.
Novelty handling is therefore not only a robustness add-on; it is part of the basic problem of autonomy.

The consequence is that novelty cannot be evaluated independently of the agent.
A heavier object may be a simple parameter update for a controller that already represents mass, but an unrepresentable change for an agent whose state description has no place for that distinction.
A new object may be immediately meaningful to an agent with the right category and useless to one without it.
The structure of the agent's knowledge determines what can be absorbed, what must be reconstructed, and where adaptation can be localized.
The next section develops this difference as a distinction between two modes of response to mismatch: assimilation and accommodation.
